% Formal model of evidence-gated completion (Sec. IV-C), as implemented after the fixes (commit eb3803e).
\subsection{Evidence-Gated Verification Model}
\label{sec:formal}

We model an episode as a sequence of browser states $s_t$ with observations $o_t=\Omega(s_t)$ (URL, title, element manifest, screenshot). The planner maps the goal $g$, $o_t$ and retrieved context to an action $a_t$; the executor applies it and returns a dispatch flag $\alpha_t\in\{0,1\}$, which is $1$ when the browser call raised no exception and, for text entry, the value read back matched. Some fallback paths (a mouse-and-keyboard retry of text entry, vision-coordinate clicks) return $\alpha_t=1$ without such a check.

\paragraph{Blocking predicate} $\beta(s)=1$ when the page is a bot check: the URL contains a challenge path, the title contains challenge wording (\emph{captcha}, \emph{not a robot}, \emph{sorry\ldots}), the body contains one of 12 challenge phrases, or a visible challenge frame other than an invisible reCAPTCHA is present. This check runs before $\Phi$ and ends the episode as \textsc{Blocked}.

\paragraph{Step and task predicates} The decomposer produces a plan $P=(g_1,\dots,g_n)$. The per-step predicate is still weak, $\phi_i := \alpha_t$: a step is marked complete when an action other than \texttt{complete} dispatches during it. Completion is decided by $\Phi(g,s_{t+1},\hat{y},P)$, evaluated on the live page after the action:
\begin{equation}
\Phi = \neg E \wedge \neg H \wedge Q \wedge T \wedge X \wedge R,
\label{eq:phi}
\end{equation}
where $E$ holds for a browser-error URL, an error or ``not found'' title, or a redirect to a sign-in page that the request did not name; $H$ holds when a search or extraction request, or a multi-step plan, is still on the bare home page of one of four search engines; $Q$ requires a results page to carry the query quoted in $g$; $T$ requires the literal text that $g$ asks to be entered to appear in an input, text area or editable element, exactly if $g$ says ``exactly''; $X$, for extraction requests, requires an answer or structured extraction that mentions at least one content word of $g$ and does not merely restate a navigation step; and $R$, for multi-step plans, requires that no step is pending unless a structured extraction exists. $H$, $Q$, $T$, $X$ and $R$ are instantiated only when their triggers occur; when none does, $\Phi$ reduces to $\neg E$.

\paragraph{When $\Phi$ is evaluated} $\Phi$ is evaluated when the planner emits \texttt{complete}, when the plan index has reached $n$, or when a requested text entry has dispatched. For a one-step plan, the default for short requests, this is after every successful action, so $\Phi$ also acts as the termination rule. The answer $\hat{y}$ is the text an extraction grounded in page lines, or else the planner's summary on \texttt{complete}; the reasoning attached to ordinary actions is not an answer.

\paragraph{Completion rule} The agent reports success only if $\Phi=1$; each evaluation is written to the evidence store with its expected and observed values. A passing evaluation can still yield \textsc{Waiting approval}, and the terminal node downgrades a success to \textsc{Failed} if the intent named a site but no navigation was verified. A rejected evaluation does not end the episode: control returns to observation and planning until a counter of intent, decomposition and planning calls reaches $K=15$.

\paragraph{Consequence} Dispatch success does not imply reported success, $\alpha_t = 1 \not\Rightarrow \Phi=1$; a successful \texttt{navigate} to a search engine's home page for a search request is an example ($H=1$). The converse, that $\Phi=1$ implies the goal is met, does not hold, and Section~\ref{sec:results} measures how often it fails.

\paragraph{False completion rate} Let $\mathcal{T}_c$ be the episodes reported complete and $\Phi^{\star}$ an independent, task-specific oracle written by the experimenter. Then
\begin{equation}
\mathrm{FCR} = \frac{\lvert\{\tau\in\mathcal{T}_c : \Phi^{\star}(\tau)=0\}\rvert}{\lvert\mathcal{T}_c\rvert}.
\label{eq:fcr}
\end{equation}
FCR is a precision-type measure; an agent that never completes has no false completions. It must therefore be reported together with task success and the false-rejection rate, the fraction of goal-satisfying episodes not reported complete.
